% SPDX-License-Identifier: GPL-3.0-or-later
% Copyright (C) 2026 Alessandro Greco
%% 简体中文版 EPUB 替代文本。
\altfig{Tor 的五项密码学任务：保密性由信元所用 AES-CTR 及 TLS 通道所用 AES-GCM 或 ChaCha20-Poly1305 实现（第 10、7 章）；完整性涉及 SHA-1、SHA3-256 摘要和 CGO 中基于 POLYVAL 的认证标签（第 10 章）；ntor 与 ntor-v3 使用 X25519 协商密钥（第 9 章）；证书和共识使用 Ed25519 及旧式 RSA 认证（第 7、12 章）；电路密钥通过 HKDF-SHA256 和 SHAKE-256 派生（第 9 章）。}
% Tor 中的五类密码学任务及其原语。
\begin{tikzpicture}[font=\sffamily\scriptsize,
  compito/.style={draw=torviola,fill=torviola!8,rounded corners=2pt,minimum width=2.6cm,minimum height=0.7cm},
  prim/.style={draw=torverde!60!black,fill=torverde!12,rounded corners=2pt,text width=3.9cm,align=center,minimum height=0.7cm},
  dove/.style={font=\sffamily\tiny,text=torgrigio,anchor=west}]
  \foreach \i/\c/\p/\d in {%
    0/机密性/{AES-CTR；TLS 使用 AES-GCM 或 ChaCha20-Poly1305}/{信元（第 10 章），通道（第 7 章）},
    1/完整性/{SHA-1、SHA3-256、POLYVAL}/{RELAY 信元摘要、CGO 认证标签（第 10 章）},
    2/{密钥协商}/{X25519 (Diffie-Hellman)}/{ntor、ntor-v3（第 9 章）},
    3/身份认证/{Ed25519、RSA（逐步淘汰）}/{证书、共识（第 7、12 章）},
    4/密钥派生/{HKDF-SHA256、SHAKE-256}/{电路密钥（第 9 章）}} {
    \node[compito] (c\i) at (0,-\i*1.0) {\c};
    \node[prim] (p\i) at (4.4,-\i*1.0) {\p};
    \draw[->,torgrigio] (c\i) -- (p\i);
    \node[dove] at (6.8,-\i*1.0) {\d};
  }
\end{tikzpicture}
